\section{Including MATLAB or Python Code}

This template includes a pre-configured \texttt{matlabstyle} for \texttt{lstlisting} that provides syntax highlighting and a nice background box for your code.

To include code directly in your document, use the \texttt{lstlisting} environment and set the style to \texttt{matlabstyle}.

\begin{lstlisting}[style=matlabstyle, caption={A simple MATLAB script to plot a sine wave.}]
% Define the time vector
t = 0:0.01:2*pi;

% Generate the signal
x = sin(2*pi*1*t);

% Plot the signal
figure;
plot(t, x, 'LineWidth', 1.5);
title('Sine Wave');
xlabel('Time (s)');
ylabel('Amplitude');
grid on;
\end{lstlisting}

Even though there's a \texttt{matlabstyle}, you can also use the newly added \texttt{pythonstyle} to display Python code cleanly with appropriate syntax highlighting:

\begin{lstlisting}[style=pythonstyle, caption={A simple Python script.}]
import numpy as np
import matplotlib.pyplot as plt

# Define the time vector
t = np.linspace(0, 2*np.pi, 100)
x = np.sin(2 * np.pi * 1 * t)

plt.plot(t, x)
plt.title('Sine Wave')
plt.show()
\end{lstlisting}

We also provide \texttt{cppstyle} for C/C++ developers:

\begin{lstlisting}[style=cppstyle, caption={A simple C++ program.}]
#include <iostream>
#include <vector>

int main() {
    std::vector<int> numbers = {1, 2, 3, 4, 5};
    for (int num : numbers) {
        std::cout << "Number: " << num << std::endl;
    }
    return 0;
}
\end{lstlisting}

For web developers, there is \texttt{jsstyle} which supports JavaScript and TypeScript natively:

\begin{lstlisting}[style=jsstyle, caption={A modern TypeScript function.}]
interface User {
    id: number;
    name: string;
}

const fetchUser = async (id: number): Promise<User> => {
    const response = await fetch(`/api/users/${id}`);
    const data = await response.json();
    return data;
};
\end{lstlisting}

Finally, you can even showcase LaTeX code itself using the \texttt{latexstyle}:

\begin{lstlisting}[style=latexstyle, caption={A snippet of LaTeX code.}]
\begin{problem}[1: Signal Analysis]
Calculate the energy and power of the given signal $x(t) = e^{-2t}u(t)$.
\end{problem}
\end{lstlisting}

\newpage